\chapter*{Sommario}
\addcontentsline{toc}{chapter}{Sommario}

I pesi dei modelli linguistici Transformer determinano gran parte della memoria richiesta e del traffico di dati durante la generazione del testo, soprattutto sui dispositivi personali, dove la memoria è la risorsa più scarsa. Questa tesi propone STERA, una parametrizzazione strutturata delle proiezioni lineari. Ogni proiezione è composta da blocchi $128\times128$ di tipo RowMask23, ciascuno descritto da 3.176 coordinate anziché da 16.384 coefficienti: una diagonale libera, una base condivisa lungo le diagonali generalizzate dello XOR degli indici e 23 correzioni per riga. I blocchi sono completati da una scala per blocco, da guadagni per riga e per colonna e da un termine a basso rango. La mappa dalle coordinate al blocco è lineare e iniettiva, il gradiente rispetto alle coordinate ha forma chiusa e il prodotto con un vettore si calcola con la trasformata di Walsh--Hadamard, senza ricostruire la matrice.

La parametrizzazione è applicata a STERA 30M, un Transformer a sei strati con 30,3 milioni di parametri allenabili, addestrato da zero sul corpus TinyStories per 1,53 miliardi di token. Su 300 racconti di prova il modello ottiene una NLL di 1,174 nat per token, inferiore a quella dei checkpoint pubblici TinyStories-8M (1,515) e TinyStories-33M (1,324). In una valutazione cieca delle continuazioni generate per gli stessi incipit, il 44,0\% dei racconti di STERA 30M è giudicato corretto, contro il 21,7\% di TinyStories-8M e il 48,3\% di TinyStories-33M, una differenza non significativa rispetto al secondo, che ha più del doppio dei parametri. I blocchi RowMask23 garantiscono per costruzione l'80,6\% di valori in meno sulla parte strutturata delle matrici; con scale, guadagni e basso rango, i pesi di STERA 30M occupano 121 MB in FP32, il 63\% in meno della loro espansione densa. Applicando le proiezioni direttamente alle coordinate, il tempo per token in generazione si riduce del 34,8\% rispetto alle matrici dense equivalenti. Il termine a basso rango si quantizza a 8 bit con una variazione della NLL inferiore allo 0,03\% e una quota di racconti corretti non significativamente diversa, portando i pesi a 86 MB e rendendo la generazione più rapida.

\vspace{1.5em}
\section*{Abstract}

The weights of Transformer language models account for most of the memory and data traffic required during text generation, especially on personal devices, where memory is the scarcest resource. This thesis proposes STERA, a structured parameterization of linear projections. Each projection is composed of $128\times128$ RowMask23 blocks, each described by 3,176 coordinates instead of 16,384 coefficients: a free diagonal, a basis shared along the generalized diagonals defined by the XOR of the indices, and 23 corrections per row. The blocks are complemented by one scale per block, row and column gains, and a low-rank term. The map from coordinates to block is linear and injective, the gradient with respect to the coordinates has a closed form, and the matrix--vector product is computed with the Walsh--Hadamard transform without reconstructing the matrix.

The parameterization is applied to STERA 30M, a six-layer Transformer with 30.3 million trainable parameters, trained from scratch on the TinyStories corpus for 1.53 billion tokens. On 300 test stories the model reaches a negative log-likelihood of 1.174 nats per token, lower than the public TinyStories-8M (1.515) and TinyStories-33M (1.324) checkpoints. In a blind evaluation of continuations generated from the same prompts, 44.0\% of STERA 30M stories are judged correct, against 21.7\% for TinyStories-8M and 48.3\% for TinyStories-33M, a non-significant difference with respect to the latter, which has more than twice as many parameters. The RowMask23 blocks guarantee by construction 80.6\% fewer values in the structured part of the matrices; with scales, gains and the low-rank term, the weights of STERA 30M take 121 MB in FP32, 63\% less than their dense expansion. Applying the projections directly to the coordinates reduces the per-token generation time by 34.8\% compared with the equivalent dense matrices. The low-rank term can be quantized to 8 bits with a change in negative log-likelihood below 0.03\% and a non-significantly different share of correct stories, bringing the weights to 86 MB and making generation faster.
